% F085｜真实数学图；代表函数不构成抽象解的唯一性证明。
\begin{center}
\begin{minipage}{.49\linewidth}\centering
\begin{tikzpicture}
\begin{axis}[width=6.1cm,height=4.8cm,axis lines=middle,
 xmin=-2.9,xmax=2.9,ymin=-1.5,ymax=3.4,
 xtick={-2,0,2},xticklabels=\empty,ytick={-1,0,1,2,3},xlabel={$t$},ylabel={$f(t)$},
 title={同属余弦型，形状却不同},title style={font=\small,yshift=3pt},
 tick label style={font=\small},label style={font=\small},
 ylabel style={at={(axis description cs:0,1)},anchor=south east,xshift=-18pt},
 clip=false,grid=major,grid style={black!8},enlargelimits=false,
 every axis plot/.append style={thick},samples=161]
\addplot[CExplanation,very thick,domain=-2.5:2.5] {cos(deg(x))};
\addplot[orange!80!black,dashed,domain=-2.5:2.5] {(2^x+2^(-x))/2};\addplot[green!45!black,dotted,domain=-2.5:2.5] {1};
\node[font=\small,anchor=north,yshift=-3pt] at (axis cs:-2,-1.5) {$-2$};
\node[font=\small,anchor=north,yshift=-3pt] at (axis cs:2,-1.5) {$2$};

\end{axis}
\end{tikzpicture}
\end{minipage}\hfill
\begin{minipage}{.49\linewidth}\centering
\begin{tikzpicture}
\begin{axis}[width=6.1cm,height=4.8cm,axis lines=middle,
 xmin=-2.9,xmax=2.9,ymin=-3.1,ymax=3.1,
 xtick={-2,0,2},xticklabels=\empty,ytick={-2,0,2},xlabel={$t$},ylabel={$f(t)$},
 title={同属正弦型：曲线与直线},title style={font=\small,yshift=3pt},
 tick label style={font=\small},label style={font=\small},
 ylabel style={at={(axis description cs:0,1)},anchor=south east,xshift=-18pt},
 clip=false,grid=major,grid style={black!8},enlargelimits=false,
 every axis plot/.append style={thick},samples=161]
\addplot[CExplanation,very thick,domain=-2.5:2.5] {sin(deg(x))};
\addplot[orange!80!black,dashed,domain=-2.5:2.5] {x};
\node[font=\small,anchor=north,yshift=-3pt] at (axis cs:-2,-3.1) {$-2$};
\node[font=\small,anchor=north,yshift=-3pt] at (axis cs:2,-3.1) {$2$};

\end{axis}
\end{tikzpicture}
\end{minipage}
\par\small 左图：蓝实线 $\cos t$，橙虚线 $(2^t+2^{-t})/2$，绿点线 $1$；右图：蓝实线 $\sin t$，橙虚线 $t$。同一模型容纳不同曲线，不能直接断言周期、值域或唯一解析式。
\end{center}
